% Chapter II, Part 1: The Benard Problem (Sections 5-15)
% PDF pages 25-58, book pages 9-42

\chapter{THE THERMAL INSTABILITY OF A LAYER OF FLUID HEATED FROM BELOW}
\label{ch:2}

\begin{center}
  {\Large 1. THE B\'ENARD PROBLEM}
\end{center}

\section{Introduction}\label{sec:5}

The problem of the onset of thermal instability in horizontal layers of
fluid heated from below is well suited to illustrate the many facets,
mathematical and physical, of the general theory of hydrodynamic
stability. If we enlarge the problem to include the effects of rotation
and magnetic field, we can also exhibit the conflicting tendencies to
which a fluid can be subject; and many striking aspects of fluid behaviour,
disclosed by the mathematical theory in these connexions, have been
experimentally demonstrated. For these reasons, we have selected the
general topic of the instability of layers of fluid heated from below for a
detailed treatment; and Chapters II--V are devoted to this problem.

In this chapter we shall consider the simplest problem in the absence
of rotation and magnetic field; the effects of these will be considered in
the following chapters.

\section{The nature of the physical problem}\label{sec:6}

Consider then a horizontal layer of fluid in which an adverse temperature
gradient is maintained by heating the underside. The temperature
gradient thus maintained is qualified as adverse since, on account of
thermal expansion, the fluid at the bottom will be lighter than the fluid
at the top; and this is a top-heavy arrangement which is potentially
unstable. Because of this latter instability there will be a tendency on
the part of the fluid to redistribute itself and remedy the weakness in
its arrangement. However, this natural tendency on the part of the
fluid will be inhibited by its own viscosity. In other words, we expect
that the adverse temperature gradient which is maintained must exceed
a certain value before the instability can manifest itself.

The earliest experiments to demonstrate in a definitive manner the
onset of thermal instability in fluids are those of B\'enard in 1900, though
the phenomenon of thermal convection itself had been recognized earlier
by Count Rumford (1797) and James Thomson (1882). In \S~18 we shall
describe the experiments of B\'enard and others. It will suffice to summarize here the principal facts established by them. They are: \emph{first},
a certain critical adverse temperature gradient must be exceeded before
instability can set in; and \emph{second}, the motions which ensue on surpassing
the critical temperature gradient have a stationary cellular character.
What actually happens at the onset of instability is that the layer of
fluid resolves itself into a number of cells; and if the experiment is
performed with sufficient care, the cells become equal and they align
themselves to form a regular hexagonal pattern. Fig.~1 (which is a
reproduction of one of B\'enard's early photographs) illustrates the
phenomenon.

\figurePlaceholder{1}{B\'enard cells in spermaceti. A reproduction of one of B\'enard's original photographs.}

The theoretical foundations for a correct interpretation of the foregoing facts were laid by Lord Rayleigh in a fundamental paper. Rayleigh
showed that what decides the stability, or otherwise, of a layer of fluid
heated from below is the numerical value of the non-dimensional parameter,
\begin{equation}
\label{eq:2-1}
R = \frac{g\alpha\beta}{\kappa\nu}\,d^4,
\end{equation}
where $g$ denotes the acceleration due to gravity, $d$ the depth of the layer,
$\beta$ $(= |dT/dz|)$ the uniform adverse temperature gradient which is
maintained, and $\alpha$, $\kappa$, and $\nu$ are the coefficients of volume expansion,
thermometric conductivity, and kinematic viscosity, respectively; $R$ as
defined here is called the \emph{Rayleigh number}. Rayleigh further showed that
instability must set in when $R$ exceeds a certain critical value $R_c$, and
that when $R$ just exceeds $R_c$ a stationary pattern of motions must come
to prevail. The principal theoretical question is, therefore: how is one to
determine $R_c$? This chapter is largely devoted to answering this question.

\section{The basic hydrodynamic equations}\label{sec:7}

For a general treatment of the problem of thermal instability described
in \S~6, we need the equations governing the hydrodynamical flow of a
viscous fluid of varying density and temperature. In this section we shall
obtain these basic equations.

Consider then a fluid in which the density $\rho$ is a function of position
$x_j$ $(j = 1, 2, 3)$. Let $u_j$ $(j = 1, 2, 3)$ denote the components of the
velocity. In writing the various equations we shall use the notation
of Cartesian tensors with the usual summation convention.

\subsection{The equation of continuity}\label{sec:7a}

We have
\begin{equation}
\label{eq:2-2}
\frac{\partial\rho}{\partial t} + \frac{\partial}{\partial x_j}(\rho u_j) = 0.
\end{equation}
This equation expresses the \emph{conservation of mass}: for, by integrating the
equation over an arbitrary volume $V$, we obtain
\begin{equation}
\label{eq:2-3}
\frac{\partial}{\partial t}\int\limits_V \rho\,d\tau = -\int\limits_V \frac{\partial}{\partial x_j}(\rho u_j)\,d\tau,
\end{equation}
where $d\tau$ $(= dx_1\,dx_2\,dx_3)$ is the element of volume. By Gauss's theorem,
the volume integral on the right-hand side is equal to an integral over
the bounding surface $S$ of $V$; thus,
\begin{equation}
\label{eq:2-4}
\frac{\partial}{\partial t}\int\limits_V \rho\,d\tau = -\int\limits_S \rho u_j\,dS_j,
\end{equation}
where $dS_j$ represents a vector normal to the element of surface $dS$ and
whose absolute magnitude is equal to $dS$. [The vector $dS_j$ is dual to
the antisymmetric tensor $dS_{lm}$ whose components are equal to the projected areas of $dS$ on the coordinate planes.] Equation~(4) expresses the
fact that the rate of change of the mass contained in a fixed volume of
the fluid is given by the rate at which the fluid flows out of it across the
boundary $S$. We could, indeed, have started with equation~(4) as the
expression of the conservation of mass and retraced the arguments to
derive equation~(\ref{eq:2-2}).

An alternative form of equation~(\ref{eq:2-2}) which we shall find useful is
\begin{equation}
\label{eq:2-5}
\frac{\partial\rho}{\partial t} + u_j\frac{\partial\rho}{\partial x_j} = -\rho\frac{\partial u_j}{\partial x_j}.
\end{equation}
The quantity on the left-hand side of equation~(\ref{eq:2-5}) is clearly the total
derivative of $\rho$ with respect to the time.

For an incompressible fluid, the equation of continuity reduces to
\begin{equation}
\label{eq:2-6}
\frac{\partial u_j}{\partial x_j} = 0;
\end{equation}
in this case, the velocity field is, therefore, \emph{solenoidal}.

\subsection{The equations of motion}\label{sec:7b}

Before we can write down the equations of motion, it is necessary to
know the \emph{stress} $P_{ij}$ acting in the direction of $x_j$ per unit area on an
element of surface normal to $x_i$. This stress must be related to the \emph{rate
of increase of strain} in the fluid; the latter is given by
\begin{equation}
\label{eq:2-7}
e_{ij} = \frac{1}{2}\!\left(\frac{\partial u_i}{\partial x_j}+\frac{\partial u_j}{\partial x_i}\right)\!.
\end{equation}
A basic assumption of fluid dynamics is that $P_{ij}$ and $e_{ij}$ are related \emph{linearly};
thus,
\begin{equation}
\label{eq:2-8}
P_{ij} = \varpi_{ij} + \eta_{ij;kl}\,e_{kl},
\end{equation}
where $\varpi_{ij}$ is a symmetric tensor to which $P_{ij}$ tends in the limit $e_{ij} = 0$,
and $\eta_{ij;kl}$ is a tensor of the fourth rank. For an isotropic fluid, the form
of the relation~(\ref{eq:2-8}) must be invariant to arbitrary rotations and translations of the coordinate system. This requires that $\varpi_{ij}$ and $\eta_{ij;kl}$ are
\emph{isotropic tensors}; they must, therefore, have the forms
\begin{equation}
\label{eq:2-9}
\varpi_{ij} = -p\,\delta_{ij}
\end{equation}
and
\begin{equation}
\label{eq:2-10}
\eta_{ij;kl} = \lambda\,\delta_{ij}\,\delta_{kl} + \mu(\delta_{ik}\,\delta_{jl} + \delta_{il}\,\delta_{jk}),
\end{equation}
where $p$, $\lambda$, and $\mu$ are (arbitrary) scalar functions of $x_i$. Inserting the
expressions~(\ref{eq:2-9}) and~(\ref{eq:2-10}) in equation~(\ref{eq:2-8}), we have
\begin{equation}
\label{eq:2-11}
P_{ij} = -p\,\delta_{ij} + 2\mu\,e_{ij} + \lambda\,\delta_{ij}\,e_{kk}.
\end{equation}
It is convenient to define $p$ as the isotropic pressure at $x_i$ when there is
no strain; then,
\begin{equation}
\label{eq:2-12}
P_{kk} = -3p + 2\mu\,e_{kk} + 3\lambda\,e_{kk}.
\end{equation}
Hence, with the definition of $p$ adopted,
\begin{equation}
\label{eq:2-13}
\lambda = -\tfrac{2}{3}\mu
\end{equation}
and
\begin{equation}
\label{eq:2-14}
P_{ij} = -p\,\delta_{ij} + 2\mu\,e_{ij} - \tfrac{2}{3}\mu\,\delta_{ij}\,e_{kk}.
\end{equation}
The coefficient $\mu$ which occurs in this equation is the \emph{coefficient of viscosity};
it can be an arbitrary function of position.

The terms in~(\ref{eq:2-14}) which are proportional to $\mu$ define the stresses due
to viscosity; denoting this by $p_{ij}$, we have
\begin{equation}
\label{eq:2-15}
p_{ij} = \mu\!\left(\frac{\partial u_i}{\partial x_j}+\frac{\partial u_j}{\partial x_i}\right) - \tfrac{2}{3}\mu\,\delta_{ij}\,\frac{\partial u_k}{\partial x_k}.
\end{equation}
For an incompressible fluid, the viscous tensor has the simpler form
(cf.\ equation~(\ref{eq:2-8}))
\begin{equation}
\label{eq:2-16}
p_{ij} = \mu\!\left(\frac{\partial u_i}{\partial x_j}+\frac{\partial u_j}{\partial x_i}\right)\!.
\end{equation}

In terms of the stresses $P_{ij}$ we can write down the hydrodynamical
equations of motion. We have
\begin{equation}
\label{eq:2-17}
\rho\frac{\partial u_i}{\partial t} + \rho u_j\frac{\partial u_i}{\partial x_j} = \rho X_i + \frac{\partial P_{ij}}{\partial x_j},
\end{equation}
where $X_i$ is the $i$th component of whatever external force may be acting
on the fluid. Substituting for $P_{ij}$ in this equation, we have
\begin{equation}
\label{eq:2-18}
\rho\frac{\partial u_i}{\partial t} + \rho u_j\frac{\partial u_i}{\partial x_j} = \rho X_i - \frac{\partial p}{\partial x_i} + \frac{\partial}{\partial x_j}\!\left[\mu\!\left(\frac{\partial u_i}{\partial x_j}+\frac{\partial u_j}{\partial x_i}\right)\right] - \tfrac{2}{3}\frac{\partial}{\partial x_i}\!\left(\mu\frac{\partial u_k}{\partial x_k}\right)\!.
\end{equation}
For an incompressible fluid in which $\mu$ is constant, equation~(\ref{eq:2-18}) simplifies to
\begin{equation}
\label{eq:2-19}
\rho\frac{\partial u_i}{\partial t} + \rho u_j\frac{\partial u_i}{\partial x_j} = \rho X_i - \frac{\partial p}{\partial x_i} + \mu\nabla^2 u_i.
\end{equation}
This is the original form of the \emph{Navier--Stokes equations}.

Equation~(\ref{eq:2-17}) expresses the \emph{conservation of momentum}. This can be
seen by integrating the equation over an arbitrary volume $V$. We find
\begin{equation}
\label{eq:2-20}
\int\limits_V\!\left(\rho\frac{\partial u_i}{\partial t} + \rho u_j\frac{\partial u_i}{\partial x_j}\right)d\tau = \int\limits_V \rho X_i\,d\tau + \int\limits_V \frac{\partial P_{ij}}{\partial x_j}\,d\tau.
\end{equation}
By integrating by parts the second integral on the left-hand side and
transforming the integral over $P_{ij}$ on the right-hand side by Gauss's
theorem, we obtain
\begin{equation}
\label{eq:2-21}
\int\limits_V\!\left(\rho\frac{\partial u_i}{\partial t} - u_i\frac{\partial}{\partial x_j}(\rho u_j)\right)d\tau + \int\limits_S \rho u_i\,u_j\,dS_j = \int\limits_V \rho X_i\,d\tau + \int\limits_S P_{ij}\,dS_j.
\end{equation}
By making use of the equation of continuity, we find that the first of the
two integrals on the left-hand side of~(\ref{eq:2-21}) is
\begin{equation}
\label{eq:2-22}
\int\limits_V\!\left(\rho\frac{\partial u_i}{\partial t} + u_i\frac{\partial\rho}{\partial t}\right)d\tau = \frac{\partial}{\partial t}\int\limits_V \rho u_i\,d\tau.
\end{equation}
Equation~(\ref{eq:2-21}) can accordingly be written as
\begin{equation}
\label{eq:2-23}
\frac{\partial}{\partial t}\int\limits_V \rho u_i\,d\tau = \int\limits_V \rho X_i\,d\tau + \int\limits_S P_{ij}\,dS_j - \int\limits_S \rho u_i\,u_j\,dS_j.
\end{equation}
This equation expresses the fact that the rate of change of the momentum
contained in a fixed volume $V$ of the fluid is equal to the volume integral
of the external forces acting on the elements of the fluid plus the surface integral of the normal stresses acting on the bounding surface $S$ of $V$
minus the rate at which momentum flows out of $V$ across the boundaries
of $V$ by the motions prevailing on $S$. It is clear that by expressing the
law of the conservation of momentum in the form~(\ref{eq:2-23}), we can, conversely, derive the equation of motion~(\ref{eq:2-17}).

\subsection{The rate of viscous dissipation}\label{sec:7c}

By multiplying equation~(\ref{eq:2-17}) by $u_i$ (and, of course, summing over $i$)
and integrating over a volume $V$, we obtain
\begin{equation}
\label{eq:2-24}
\frac{1}{2}\int\limits_V \rho\frac{\partial}{\partial t}u_i^2\,d\tau + \frac{1}{2}\int\limits_V \rho u_j\frac{\partial}{\partial x_j}u_i^2\,d\tau = \int\limits_V \rho u_i\,X_i\,d\tau + \int\limits_V u_i\frac{\partial P_{ij}}{\partial x_j}\,d\tau.
\end{equation}
By a sequence of transformations similar to those used in passing from
equation~(\ref{eq:2-20}) to equation~(\ref{eq:2-23}), we find
\begin{equation}
\label{eq:2-25}
\frac{1}{2}\frac{\partial}{\partial t}\int\limits_V \rho u_i^2\,d\tau = \int\limits_V \rho u_i\,X_i\,d\tau + \int\limits_S u_i\,P_{ij}\,dS_j - \frac{1}{2}\int\limits_S \rho u_i^2\,u_j\,dS_j - \int\limits_V P_{ij}\frac{\partial u_i}{\partial x_j}\,d\tau.
\end{equation}
This equation gives the rate at which the kinetic energy contained in a
volume $V$ of the fluid changes; and we observe that this is the sum of
four terms. The first three terms represent, respectively, the rate at
which the external forces do work on the elements of the fluid in $V$; the
rate at which the stresses $P_{ij}$ do work on the surface bounding $V$; and
finally, the rate at which energy actually flows out of $V$ across $S$.

It remains to interpret the last term on the right-hand side of equation~(\ref{eq:2-25}); this is an integral over the volume of the quantity
\begin{equation}
\label{eq:2-26}
-P_{ij}\frac{\partial u_i}{\partial x_j}.
\end{equation}
Using the definitions of $e_{ij}$ and $P_{ij}$ (equations~(\ref{eq:2-7}) and~(\ref{eq:2-14})), we can write
\begin{equation}
\label{eq:2-27}
P_{ij}\frac{\partial u_i}{\partial x_j} = P_{ij}\,e_{ij} = (-p\,\delta_{ij} + 2\mu\,e_{ij} - \tfrac{2}{3}\mu\,\delta_{ij}\,e_{kk})e_{ij} = -p\,e_{jj} + 2\mu\,e_{ij}^2 - \tfrac{2}{3}\mu(e_{jj})^2.
\end{equation}
The first term on the right-hand side is
\begin{equation}
\label{eq:2-28}
-p\,e_{jj} = -p\frac{\partial u_j}{\partial x_j} = \frac{p}{\rho}\!\left(\frac{\partial\rho}{\partial t} + u_j\frac{\partial\rho}{\partial x_j}\right) = \frac{p}{\rho}\frac{d\rho}{dt};
\end{equation}
the volume integral of this quantity represents, therefore, the increase
in the internal energy due to the compression which the fluid experiences.
The remaining term,
\begin{equation}
\label{eq:2-29}
\Phi = 2\mu e_{ij}^2 - \tfrac{2}{3}\mu(e_{jj})^2,
\end{equation}
in~(\ref{eq:2-27}) must, therefore, represent the rate at which energy is dissipated,
irreversibly, by viscosity in each element of volume of the fluid. Consistent with this statement, $\Phi$ is indeed positive definite: for, it can be
readily verified that an equivalent form for $\Phi$ is
\begin{equation}
\label{eq:2-30}
\Phi = 4\mu(e_{12}^2 + e_{13}^2 + e_{23}^2) + \tfrac{4}{3}\mu[(e_{11}-e_{22})^2 + (e_{22}-e_{33})^2 + (e_{33}-e_{11})^2].
\end{equation}
For an incompressible fluid $e_{jj} = 0$, and the corresponding expression
for $\Phi$ is
\begin{equation}
\label{eq:2-31}
\Phi = 2\mu e_{ij}^2.
\end{equation}

\subsection{The equation of heat conduction}\label{sec:7d}

We have seen that the laws of conservation of mass and of momentum
lead to the equations of continuity and motion, respectively. It remains
to express the law of the \emph{conservation of energy}. As we shall presently see,
this leads to the equation of heat conduction in the form required for
our present purposes.

Apart from an additive constant, the energy $\epsilon$ per unit mass of the
fluid can be written as
\begin{equation}
\label{eq:2-32}
\epsilon = \tfrac{1}{2}u_i^2 + c_v\,T,
\end{equation}
where $c_v$ is the specific heat at constant volume and $T$ is the temperature.
Counting the gains and losses of energy which occur in a volume $V$ of
the fluid, per unit time, we have
\begin{align}
\label{eq:2-33}
\frac{\partial}{\partial t}\int\limits_V \rho\epsilon\,d\tau &= \text{rate at which work is done on the boundary } S \text{ of } V \text{ by the} \notag\\
& \quad\text{stresses } P_{ij} + {} \notag\\
& \quad{} + \text{rate at which work is done on each element of the fluid} \notag\\
& \quad\text{inside } V \text{ by the external forces---} \notag\\
& \quad{} - \text{rate at which energy in the form of heat is \emph{conducted} across } S\text{---} \notag\\
& \quad{} - \text{rate at which energy is \emph{convected} across } S \text{ by the prevailing} \notag\\
& \quad\text{mass motions} \notag\\[4pt]
&= \int\limits_S u_i\,P_{ij}\,dS_j + \int\limits_V \rho u_i\,X_i\,d\tau + \int\limits_S k\frac{\partial T}{\partial x_j}\,dS_j - \int\limits_S \rho\epsilon\,u_j\,dS_j,
\end{align}
where $k$ is the coefficient of heat conduction. Making use of equations~(\ref{eq:2-25}),~(\ref{eq:2-27}),~(\ref{eq:2-28}), and~(\ref{eq:2-29}), we can rewrite the first term on the
right-hand side of equation~(\ref{eq:2-33}) as
\begin{equation}
\label{eq:2-34}
\int\limits_S u_i\,P_{ij}\,dS_j = \frac{1}{2}\frac{\partial}{\partial t}\int\limits_V \rho u_i^2\,d\tau + \frac{1}{2}\int\limits_S \rho u_i^2\,u_j\,dS_j - \int\limits_V \rho u_i\,X_i\,d\tau - \int\limits_V p\frac{\partial u_j}{\partial x_j}\,d\tau + \int\limits_V \Phi\,d\tau;
\end{equation}
also, alternative forms of the third and the fourth terms are
\begin{equation}
\label{eq:2-35}
\int\limits_S k\frac{\partial T}{\partial x_j}\,dS_j = \int\limits_V \frac{\partial}{\partial x_j}\!\left(k\frac{\partial T}{\partial x_j}\right)d\tau
\end{equation}
and
\begin{align}
\label{eq:2-36}
-\int\limits_S \rho\epsilon\,u_j\,dS_j &= -\int\limits_V \rho(\tfrac{1}{2}u_i^2 + c_v\,T)u_j\,dS_j \notag\\
&= -\frac{1}{2}\int\limits_S \rho u_i^2\,u_j\,dS_j - \int\limits_V \frac{\partial}{\partial x_j}(\rho c_v\,T\,u_j)\,d\tau.
\end{align}
Combining the foregoing equations, we have
\begin{equation}
\label{eq:2-37}
\int\limits_V \frac{\partial}{\partial t}(\rho c_v\,T)\,d\tau = \int\limits_V \frac{\partial}{\partial x_j}\!\left(k\frac{\partial T}{\partial x_j}\right)d\tau - \int\limits_V p\frac{\partial u_j}{\partial x_j}\,d\tau + \int\limits_V \Phi\,d\tau - \int\limits_V \frac{\partial}{\partial x_j}(\rho c_v\,T\,u_j)\,d\tau.
\end{equation}
Since this equation must hold for every volume $V$, we must have
\begin{equation}
\label{eq:2-38}
\frac{\partial}{\partial t}(\rho c_v\,T) + \frac{\partial}{\partial x_j}(\rho c_v\,T\,u_j) = \frac{\partial}{\partial x_j}\!\left(k\frac{\partial T}{\partial x_j}\right) - p\frac{\partial u_j}{\partial x_j} + \Phi.
\end{equation}
Making use of the equation of continuity, we can simplify the foregoing
equation to the form
\begin{equation}
\label{eq:2-39}
\rho\frac{\partial}{\partial t}(c_v\,T) + \rho u_j\frac{\partial}{\partial x_j}(c_v\,T) = \frac{\partial}{\partial x_j}\!\left(k\frac{\partial T}{\partial x_j}\right) - p\frac{\partial u_j}{\partial x_j} + \Phi.
\end{equation}

Equations~(\ref{eq:2-2}),~(\ref{eq:2-14}),~(\ref{eq:2-17}),~(\ref{eq:2-29}), and~(\ref{eq:2-39}) are the basic hydrodynamic
equations. They must be supplemented by an equation of state. For
substances with which we shall be principally concerned, we can write
\begin{equation}
\label{eq:2-40}
\rho = \rho_0[1 - \alpha(T - T_0)],
\end{equation}
where $\alpha$ is the coefficient of volume expansion and $T_0$ is the temperature
at which $\rho = \rho_0$.

\section{The Boussinesq approximation}\label{sec:8}

In deriving the hydrodynamical equations in the preceding section,
no assumptions were made regarding the constancy, or otherwise, of the
various coefficients ($\mu$, $c_v$, $\alpha$, and $k$) which were introduced. The equations which were derived are, therefore, of quite general validity. However, as was first pointed out by Boussinesq, there are many situations
of practical occurrence in which the basic equations can be simplified
considerably. These situations occur when the variability in the
density and in the various coefficients is due to variations in the temperature of only moderate amounts. The origin of the simplifications in
these cases is due to the smallness of the coefficient of volume expansion:
for gases and liquids such as we shall be mostly concerned with, $\alpha$ is in
the range $10^{-3}$ to $10^{-4}$. For variations in temperature not exceeding $10^\circ$
(say), the variations in the density are at most 1 per cent. The variations
in the other coefficients (consequent to variations in density of the
amounts stated) must be of the same order. Variations of this small
amount can, in general, be ignored. But there is one important exception: the variability of $\rho$ in the term $\rho X_i$ in the equation of motion
cannot be ignored; this is because the acceleration resulting from $\delta\rho\,X_i = \alpha\,\Delta T\,X_i$
(where $\Delta T$ is a measure of the variations in temperature which occur)
can be quite large; larger than, for example, the acceleration due to the
inertial term $u_j\,\partial u_i/\partial x_j$ in the equation of motion. Accordingly, we may
treat $\rho$ as a constant in all terms in the equations of motion except the
one in the external force. This is the \emph{Boussinesq approximation}. We
shall verify in \S~9 that for the particular problem of thermal instability
we need not make the Boussinesq approximation: for, we shall show that
in this case the general equations of \S~7 lead to the same perturbation
equations. Nevertheless, the equations which follow on the Boussinesq
approximation are of interest in themselves; and they provide also the
basis for further developments in the non-linear domain.

On the basis of the foregoing remarks, we first replace the equation
of continuity~(\ref{eq:2-6}) by
\begin{equation}
\label{eq:2-41}
\frac{\partial u_j}{\partial x_j} = 0;
\end{equation}
for, the terms on the left-hand side of~(\ref{eq:2-5}) are of order $\alpha$ compared with
those on the right-hand side. With this condition on $u_j$, the expression
for the viscous stress-tensor is
\begin{equation}
\label{eq:2-42}
p_{ij} = \mu\!\left(\frac{\partial u_i}{\partial x_j}+\frac{\partial u_j}{\partial x_i}\right)\!,
\end{equation}
where for the same reason we may treat $\mu$ as a constant. In the frame of
these approximations, the equation of motion~(\ref{eq:2-17}) becomes
\begin{equation}
\label{eq:2-43}
\frac{\partial u_i}{\partial t} + u_j\frac{\partial u_i}{\partial x_j} = -\frac{1}{\rho_0}\frac{\partial\,\delta p}{\partial x_i} + \left(1 + \frac{\delta\rho}{\rho_0}\right)X_i + \nu\nabla^2 u_i,
\end{equation}
where $\nu$ $(= \mu/\rho_0)$ denotes the \emph{kinematic viscosity}, $\rho_0$ is the density at some
properly chosen mean temperature $T_0$, and
\begin{equation}
\label{eq:2-44}
\delta\rho = -\rho_0\,\alpha(T - T_0).
\end{equation}

Considering next the equation of heat conduction~(\ref{eq:2-39}), we can treat
$c_v$ and $k$ as constants and take them outside the differentiation signs;
and we can ignore also the term $-p\operatorname{div}\mathbf{u}$ on the right-hand side. The
term in the viscous dissipation $\Phi$ can also be ignored: for, according to
equations~(\ref{eq:2-43}) and~(\ref{eq:2-44}), the prevailing velocities are of the order
$[\alpha\,\Delta T\,|\mathbf{X}|\,d^2]$, where $d$ is a measure of the linear dimensions of a system.
Consequently, the term in $\Phi$ is of the order
\begin{equation}
\label{eq:2-45}
\mu\alpha|\mathbf{X}|d/k
\end{equation}
relative to the term arising from the conduction of heat; and this ratio
for ordinary liquids (such as water and mercury) is $10^{-7}$ or $10^{-8}$ for
$d \sim 1$~cm and $|\mathbf{X}| \sim g$ (the acceleration due to gravity). Under these
circumstances, the equation of heat conduction~(\ref{eq:2-39}) reduces to
\begin{equation}
\label{eq:2-46}
\frac{\partial T}{\partial t} + u_j\frac{\partial T}{\partial x_j} = \kappa\nabla^2 T,
\end{equation}
where $\kappa$ $(= k/\rho_0 c_v)$ is the \emph{coefficient of thermometric conductivity}.

Equations~(\ref{eq:2-41}),~(\ref{eq:2-43}),~(\ref{eq:2-44}), and~(\ref{eq:2-46}) are the basic equations in the
Boussinesq approximation.

\section{The perturbation equations}\label{sec:9}

Consider an infinite horizontal layer of fluid in which a steady adverse
temperature gradient is maintained; further, let there be no motions.
The initial state is, therefore, one in which
\begin{equation}
\label{eq:2-47}
u_j = 0 \quad\text{and}\quad T = T(\lambda_j x_j),
\end{equation}
where $\boldsymbol{\lambda} = (0,0,1)$ is a unit vector in the direction of the vertical.

When no motions are present, the hydrodynamical equations require
only that the pressure distribution is governed by the equation (cf.\
equation~(\ref{eq:2-18}))
\begin{equation}
\label{eq:2-48}
\frac{\partial p}{\partial x_i} = \rho\,X_i = -g\rho\lambda_i,
\end{equation}
where
\begin{equation}
\label{eq:2-49}
\rho = \rho_0[1 + \alpha(T_0 - T)]
\end{equation}
and $\rho_0$ and $T_0$ are the density and the temperature at the lower boundary
($z = \lambda_j x_j = 0$, say). The temperature distribution is governed by the
equation
\begin{equation}
\label{eq:2-50}
\nabla^2 T = 0.
\end{equation}
In writing equation~(\ref{eq:2-39}), under conditions of a steady state in the form~(\ref{eq:2-50}), we have assumed that the coefficient of heat conduction $k$ is a
constant independent of $T$; this assumption is justified in our present
connexion.

The solution of equation~(\ref{eq:2-50}) appropriate to the problem on hand is
\begin{equation}
\label{eq:2-51}
T = T_0 - \beta\lambda_j x_j;
\end{equation}
$\beta$ is the adverse temperature gradient which is maintained. The corresponding distribution of the density is given by
\begin{equation}
\label{eq:2-52}
\rho = \rho_0(1 + \alpha\beta\lambda_j x_j).
\end{equation}
With this expression for $\rho$, equation~(\ref{eq:2-48}) can be integrated to give
\begin{equation}
\label{eq:2-53}
p = p_0 - g\rho_0(\lambda_i x_i + \tfrac{1}{2}\alpha\beta\lambda_i\lambda_j x_i x_j).
\end{equation}

Let the initial state described by equations~(\ref{eq:2-47}),~(\ref{eq:2-51}),~(\ref{eq:2-52}), and~(\ref{eq:2-53}) be
slightly perturbed. Let $u_i$ denote the velocity in the perturbed state,
and let the altered temperature distribution be
\begin{equation}
\label{eq:2-54}
T' = T_0 - \beta\lambda_j x_j + \theta.
\end{equation}
Finally, let $\delta p$ denote the change in the pressure distribution. We shall
now obtain the linearized form of the equations of motion which govern
this perturbed state.

We shall first consider the problem on the basis of Boussinesq's
approximation. By ignoring terms of the second and higher orders in the
perturbations, equations~(\ref{eq:2-43}) and~(\ref{eq:2-46}) give (see equation~(\ref{eq:2-62}) below)
\begin{equation}
\label{eq:2-55}
\frac{\partial u_i}{\partial t} = -\frac{\partial}{\partial x_i}\!\left(\frac{\delta p}{\rho_0}\right) + g\alpha\theta\lambda_i + \nu\nabla^2 u_i
\end{equation}
and
\begin{equation}
\label{eq:2-56}
\frac{\partial\theta}{\partial t} = \beta\lambda_j u_j + \kappa\nabla^2\theta;
\end{equation}
and the velocity field remains, of course, solenoidal:
\begin{equation}
\label{eq:2-57}
\frac{\partial u_j}{\partial x_j} = 0.
\end{equation}

It is of interest to verify that we obtain the same equations~(\ref{eq:2-55})--(\ref{eq:2-57})
from the general equations of motion without first making the Boussinesq approximation. The general equations of motion are:
\begin{equation}
\label{eq:2-58}
\frac{\partial\rho}{\partial t} + u_j\frac{\partial\rho}{\partial x_j} = -\rho\frac{\partial u_j}{\partial x_j},
\end{equation}
\begin{equation}
\label{eq:2-59}
\frac{\partial T}{\partial t} + u_j\frac{\partial T}{\partial x_j} = \kappa\nabla^2 T - \frac{p}{\rho c_v}\frac{\partial u_j}{\partial x_j} + \frac{1}{\rho c_v}\Phi,
\end{equation}
and
\begin{equation}
\label{eq:2-60}
\rho\frac{\partial u_i}{\partial t} + \rho u_j\frac{\partial u_i}{\partial x_j} = -\frac{\partial p}{\partial x_i} + \frac{\partial}{\partial x_j}\!\left[\mu\!\left(\frac{\partial u_i}{\partial x_j}+\frac{\partial u_j}{\partial x_i}\right)\right] - \tfrac{2}{3}\frac{\partial}{\partial x_i}\!\left(\mu\frac{\partial u_k}{\partial x_k}\right) - g\rho\lambda_i.
\end{equation}

In the unperturbed state
\begin{equation}
\label{eq:2-61}
\frac{\partial p}{\partial x_i} = \lambda_i\,\rho_0\,\alpha\beta;
\end{equation}
and the change in the density $\delta\rho$ caused by the perturbation $\theta$ in the
temperature is given by
\begin{equation}
\label{eq:2-62}
\delta\rho = -\alpha\rho\,\theta = -\alpha\rho_0(1 + \alpha\beta\lambda_j x_j)\theta.
\end{equation}
Consequently, while the first-order terms on the left-hand side of
equation~(\ref{eq:2-58}) occur with the factor $\alpha$, the corresponding terms on the
right-hand side do not have this factor; and since $\alpha \sim 10^{-3}$--$10^{-4}$, we
may neglect the term $\lambda_j u_j\alpha\beta$ on the left-hand side and \emph{deduce} the
solenoidal character of $\mathbf{u}$. Making use of this fact and remembering
that the dissipation term $\Phi$ is of the second order in $\mathbf{u}$ in any case, we
obtain for $\theta$ the equation
\begin{equation}
\label{eq:2-63}
\frac{\partial\theta}{\partial t} = \beta\lambda_j u_j + \kappa\nabla^2\theta;
\end{equation}
this is the same equation as~(\ref{eq:2-56}).

Considering next equation~(\ref{eq:2-60}), we first observe that we may treat $\mu$
as a constant; for, the variations in $\mu$ must be of order $\alpha\,\delta\rho$ and we may
ignore variations of this order when they occur multiplied by $u_i$. Thus
we obtain the linearized equation
\begin{equation}
\label{eq:2-64}
\frac{\partial u_i}{\partial t} = -\frac{1}{\rho}\frac{\partial}{\partial x_i}\,\delta p + \frac{\mu}{\rho}\nabla^2 u_i + g\alpha\theta\lambda_i;
\end{equation}
and in this equation we may clearly write $\rho_0$ in place of $\rho$. We thus
recover the equations which were obtained earlier on the Boussinesq
approximation.

Returning to equations~(\ref{eq:2-55})--(\ref{eq:2-57}), we eliminate the term $\partial p/\rho_0$ in~(\ref{eq:2-55}) by applying the operator,
\begin{equation}
\label{eq:2-65}
\operatorname{curl}_k = \epsilon_{ijk}\frac{\partial}{\partial x_j},
\end{equation}
to the $k$-component of the equation. Letting
\begin{equation}
\label{eq:2-66}
\omega_i = \epsilon_{ijk}\frac{\partial u_k}{\partial x_j}
\end{equation}
denote the \emph{vorticity}, we have the equation
\begin{equation}
\label{eq:2-67}
\frac{\partial\omega_i}{\partial t} = g\alpha\epsilon_{ijk}\frac{\partial\theta}{\partial x_j}\lambda_k + \nu\nabla^2\omega_i.
\end{equation}
Taking the curl of this equation once again, we have
\begin{equation}
\label{eq:2-68}
\frac{\partial}{\partial t}\frac{\partial\omega_k}{\partial x_j}\epsilon_{ijk} = g\alpha\epsilon_{ijk}\epsilon_{klm}\frac{\partial^2\theta}{\partial x_j\partial x_l}\lambda_m + \nu\nabla^2\frac{\partial\omega_k}{\partial x_j}\epsilon_{ijk}.
\end{equation}
Making use of the identity
\begin{equation}
\label{eq:2-69}
\epsilon_{ijk}\,\epsilon_{klm} = \delta_{il}\,\delta_{jm} - \delta_{im}\,\delta_{jl},
\end{equation}
we find
\begin{equation}
\label{eq:2-70}
\epsilon_{ijk}\frac{\partial\omega_k}{\partial x_j} = \epsilon_{ijk}\,\epsilon_{klm}\frac{\partial^2 u_m}{\partial x_j\partial x_l} = \frac{\partial}{\partial x_i}\!\left(\frac{\partial u_l}{\partial x_l}\right) - \nabla^2 u_i = -\nabla^2 u_i.
\end{equation}
Similarly,
\begin{equation}
\label{eq:2-71}
\epsilon_{ijk}\,\epsilon_{klm}\frac{\partial^2\theta}{\partial x_j\partial x_l}\lambda_m = \lambda_i\frac{\partial^2\theta}{\partial x_j\partial x_j} - \lambda_j\frac{\partial^2\theta}{\partial x_i\partial x_j} = \lambda_i\nabla^2\theta.
\end{equation}
Thus, equation~(\ref{eq:2-68}) becomes
\begin{equation}
\label{eq:2-72}
\frac{\partial}{\partial t}\nabla^2 u_i = g\alpha\!\left(\lambda_i\nabla^2\theta - \lambda_j\frac{\partial^2\theta}{\partial x_i\partial x_j}\right) + \nu\nabla^4 u_i.
\end{equation}

Now multiplying equations~(\ref{eq:2-67}) and~(\ref{eq:2-72}) by $\lambda_i$, we get
\begin{equation}
\label{eq:2-73}
\frac{\partial\zeta}{\partial t} = \nu\nabla^2\zeta
\end{equation}
and
\begin{equation}
\label{eq:2-74}
\frac{\partial}{\partial t}\nabla^2 w = g\alpha\!\left(\frac{\partial^2\theta}{\partial x_1^2}+\frac{\partial^2\theta}{\partial x_2^2}\right) + \nu\nabla^4 w,
\end{equation}
where
\begin{equation}
\label{eq:2-75}
\zeta = \lambda_i\omega_i \quad\text{and}\quad w = \lambda_i u_i
\end{equation}
are the $z$-components of the vorticity and the velocity. We also have
the equation (cf.\ equation~(\ref{eq:2-56}))
\begin{equation}
\label{eq:2-76}
\frac{\partial\theta}{\partial t} = \beta w + \kappa\nabla^2\theta.
\end{equation}

Equations~(\ref{eq:2-73}),~(\ref{eq:2-74}), and~(\ref{eq:2-76}) are the required perturbation equations.
We must seek solutions of these equations which satisfy certain boundary
conditions. We shall now formulate these conditions.

\subsection{The boundary conditions}\label{sec:9a}

The fluid is confined between the planes $z = 0$ and $z = d$; on these
two planes certain boundary conditions must be satisfied. Regardless
of the nature of these bounding surfaces, we must require
\begin{equation}
\label{eq:2-77}
\theta = 0 \quad\text{and}\quad w = 0 \qquad\text{for } z = 0 \text{ and } d;
\end{equation}
for, the surfaces $z = 0$ and $d$ are maintained at constant temperatures
and as such they can suffer no change; and it is also clear that the normal
component of the velocity must vanish on these surfaces. There are two
further boundary conditions which, however, depend on the nature of
the surfaces at $z = 0$ and $d$.

We shall distinguish two kinds of bounding surfaces: \emph{rigid surfaces} on
which no slip occurs and \emph{free surfaces} on which no tangential stresses act.

Consider first a rigid surface. The condition that no slip occurs on
this surface implies that not only $w$, but also the horizontal components
of the velocity, $u$ and $v$, vanish. Thus
\begin{equation}
\label{eq:2-78}
u = 0 \quad\text{and}\quad v = 0 \quad\text{in addition to}\quad w = 0 \;\text{on a rigid surface.}
\end{equation}
Since this condition must be satisfied for all $x$ and $y$ on the surface, it
follows from the equation of continuity,
\begin{equation}
\label{eq:2-79}
\frac{\partial u}{\partial x} + \frac{\partial v}{\partial y} + \frac{\partial w}{\partial z} = 0,
\end{equation}
that
\begin{equation}
\label{eq:2-80}
\frac{\partial w}{\partial z} = 0 \qquad\text{on a rigid surface.}
\end{equation}

The conditions on a free surface are
\begin{equation}
\label{eq:2-81}
P_{xz} = P_{yz} = 0.
\end{equation}
Since the isotropic term $-p\,\delta_{ij}$ has no transverse component, the condition~(\ref{eq:2-81}) is equivalent to the vanishing of the components $p_{xz}$ and $p_{yz}$ of
the viscous stress tensor:\footnote{It will be observed that we are applying the conditions~(\ref{eq:2-81}) on the undisplaced surface; in so doing we are ignoring, for example, the effects arising from the existence of gravity waves (see Chapter X, \S~94). If one should want to include such effects, then the conditions~(\ref{eq:2-81}) should be applied to the surface as displaced by the perturbation.}
\begin{equation}
\label{eq:2-82}
p_{xz} = \mu\!\left(\frac{\partial u}{\partial z}+\frac{\partial w}{\partial x}\right) \quad\text{and}\quad p_{yz} = \mu\!\left(\frac{\partial v}{\partial z}+\frac{\partial w}{\partial y}\right)\!.
\end{equation}
Since $w$ vanishes (for all $x$ and $y$) on the bounding surface, it follows from~(\ref{eq:2-82}) that
\begin{equation}
\label{eq:2-83}
\frac{\partial u}{\partial z} = \frac{\partial v}{\partial z} = 0 \qquad\text{on a free surface.}
\end{equation}
From the equation of continuity~(\ref{eq:2-79}) differentiated with respect to $z$,
we conclude that
\begin{equation}
\label{eq:2-84}
\frac{\partial^2 w}{\partial z^2} = 0 \qquad\text{on a free surface.}
\end{equation}

The boundary conditions on the normal component of the vorticity $\zeta$
can be deduced from the foregoing. Since
\begin{equation}
\label{eq:2-85}
\zeta = \frac{\partial v}{\partial x} - \frac{\partial u}{\partial y},
\end{equation}
it follows from equations~(\ref{eq:2-78}) and~(\ref{eq:2-83}) that
\begin{equation}
\label{eq:2-86}
\zeta = 0 \qquad\text{on a rigid surface}
\end{equation}
and
\begin{equation}
\label{eq:2-87}
\frac{\partial\zeta}{\partial z} = 0 \qquad\text{on a free surface.}
\end{equation}

\section{The analysis into normal modes}\label{sec:10}

According to the general ideas described in Chapter~I (\S~3), we must
analyse an arbitrary disturbance into a complete set of normal modes
and examine the stability of each of these modes, individually. For the
problem on hand, the analysis can be made in terms of two-dimensional
periodic waves of assigned wave numbers. Thus, we ascribe to all
quantities describing the perturbation a dependence on $x$, $y$, and $t$ of
the form
\begin{equation}
\label{eq:2-88}
\exp[i(k_x\,x + k_y\,y) + pt],
\end{equation}
where
\begin{equation}
\label{eq:2-89}
k = \sqrt{k_x^2 + k_y^2}
\end{equation}
is the wave number of the disturbance and $p$ is a constant (which can be
complex). As we have explained in \S~3, the solution of the stability
problem requires the specification of the states, for each $k$, which are
characterized by the real part of $p$ being zero.

In accordance with the remarks in the preceding paragraph, we
suppose that the perturbations $\theta$, $w$, and $\zeta$ have the forms:
\begin{equation}
\label{eq:2-90}
w = W(z)\exp[i(k_x\,x + k_y\,y) + pt], \quad
\theta = \Theta(z)\exp[i(k_x\,x + k_y\,y) + pt], \quad
\zeta = Z(z)\exp[i(k_x\,x + k_y\,y) + pt].
\end{equation}
For functions with this dependence on $x$, $y$, and $t$,
\begin{equation}
\label{eq:2-91}
\frac{\partial}{\partial t} = p, \quad \frac{\partial^2}{\partial x^2}+\frac{\partial^2}{\partial y^2} = -k^2, \quad\text{and}\quad \nabla^2 = \frac{d^2}{dz^2} - k^2;
\end{equation}
and equations~(\ref{eq:2-73}),~(\ref{eq:2-74}), and~(\ref{eq:2-76}) become
\begin{equation}
\label{eq:2-92}
p\!\left(\frac{d^2}{dz^2}-k^2\right)\!W = -g\alpha k^2\Theta + \nu\!\left(\frac{d^2}{dz^2}-k^2\right)^{\!2}W,
\end{equation}
\begin{equation}
\label{eq:2-93}
p\,\Theta = \beta W + \kappa\!\left(\frac{d^2}{dz^2}-k^2\right)\!\Theta,
\end{equation}
and
\begin{equation}
\label{eq:2-94}
pZ = \nu\!\left(\frac{d^2}{dz^2}-k^2\right)\!Z.
\end{equation}
Solutions of these equations must be sought which satisfy the boundary
conditions
\begin{equation}
\label{eq:2-95}
\Theta = 0 \;\text{ and }\; W = 0 \qquad\text{for } z = 0 \text{ and } d,
\end{equation}
\begin{equation}
\label{eq:2-96}
\left.\begin{aligned}
& Z = 0 \;\text{ and }\; \frac{dW}{dz} = 0 &&\text{on a rigid surface}\\
& \frac{dZ}{dz} = 0 \;\text{ and }\; \frac{d^2W}{dz^2} = 0 &&\text{on a free surface}
\end{aligned}\right\}
\end{equation}

It is convenient to discuss equations~(\ref{eq:2-92})--(\ref{eq:2-94}) in non-dimensional
variables. Choose the units
\begin{equation}
\label{eq:2-97}
[L] = d \quad\text{and}\quad [T] = d^2/\nu
\end{equation}
and let
\begin{equation}
\label{eq:2-98}
a = kd \quad\text{and}\quad \sigma = pd^2/\nu
\end{equation}
be the wave number and the `time constant' in these units. We shall,
however, let $x$, $y$, $z$ stand for the coordinates in the new unit of length $d$.
Equations~(\ref{eq:2-92}) and~(\ref{eq:2-93}) become
\begin{equation}
\label{eq:2-99}
(D^2 - a^2)(D^2 - a^2 - \sigma)W = -\left(\frac{g\alpha}{\kappa\nu}\,d^4\right)a^2\Theta
\end{equation}
and
\begin{equation}
\label{eq:2-100}
(D^2 - a^2 - p\sigma)\Theta = -\left(\frac{\beta}{\kappa}\,d^2\right)W,
\end{equation}
where $D = d/dz$ and $\mathrm{p}$ $(= \nu/\kappa)$ is the Prandtl number. [Observe that
$W$ and $\Theta$ have their usual dimensions: these are not measured in the
units which follow from~(\ref{eq:2-97}).] The associated boundary conditions are
\begin{equation}
\label{eq:2-101}
\Theta = 0, \quad W = 0 \quad\text{for } z = 0 \text{ and } 1,
\end{equation}
\begin{equation}
\label{eq:2-102}
DW = 0 \;\text{ for } z = 0 \text{ and } 1 \quad\text{if both bounding surfaces are rigid}
\end{equation}
or
\begin{equation}
\label{eq:2-103}
DW = 0 \;\text{ for } z = 0 \quad\text{and}\quad D^2W = 0 \;\text{ for } z = 1
\end{equation}
if the bottom surface is rigid and the top surface is free.

By eliminating $\Theta$ between the equations~(\ref{eq:2-99}) and~(\ref{eq:2-100}), we obtain
\begin{equation}
\label{eq:2-104}
(D^2-a^2)(D^2-a^2-\sigma)(D^2-a^2-\mathrm{p}\sigma)W = -Ra^2 W,
\end{equation}
where
\begin{equation}
\label{eq:2-105}
R = \frac{g\alpha\beta}{\kappa\nu}\,d^4
\end{equation}
is the Rayleigh number. An identical equation governs $\Theta$.

\subsection{The solutions for the horizontal components of the velocity}\label{sec:10a}

Once we have obtained the proper solutions of equations~(\ref{eq:2-99}) and~(\ref{eq:2-100}) we can complete the solution of the problem by determining the
horizontal components of the velocity as follows.

Express $u$ and $v$ in terms of two functions $\phi$ and $\psi$ in the manner
\begin{equation}
\label{eq:2-106}
u = \frac{\partial\phi}{\partial x}+\frac{\partial\psi}{\partial y} \quad\text{and}\quad v = \frac{\partial\phi}{\partial y}-\frac{\partial\psi}{\partial x}.
\end{equation}
Then
\begin{equation}
\label{eq:2-107}
\frac{\partial w}{\partial z} = \frac{\partial u}{\partial x}+\frac{\partial v}{\partial y} = \frac{\partial^2\phi}{\partial x^2}+\frac{\partial^2\psi}{\partial x\partial y}+\frac{\partial^2\phi}{\partial y^2}-\frac{\partial^2\psi}{\partial x\partial y} = -a^2\phi
\end{equation}
and
\begin{equation}
\label{eq:2-108}
d\zeta = \frac{\partial v}{\partial x}-\frac{\partial u}{\partial y} = \frac{\partial^2\phi}{\partial x\partial y}-\frac{\partial^2\psi}{\partial x^2}-\frac{\partial^2\phi}{\partial x\partial y}-\frac{\partial^2\psi}{\partial y^2} = -a^2\psi.
\end{equation}
Hence,
\begin{equation}
\label{eq:2-109}
\phi = -\frac{1}{k^2}\frac{\partial w}{\partial z} \quad\text{and}\quad \psi = -\frac{d}{k^2}\zeta.
\end{equation}
and we can write
\begin{align}
\label{eq:2-110}
u &= \frac{1}{a^2}\!\left(\frac{\partial^2 w}{\partial x\partial z}-\frac{\partial \zeta}{\partial y}\right) = \frac{1}{a^2}(a_x\,DW + a_y\,dZ)\exp[i(a_x\,x+a_y\,y)+\sigma t] \\
\label{eq:2-111}
v &= \frac{1}{a^2}\!\left(\frac{\partial^2 w}{\partial y\partial z}+\frac{\partial \zeta}{\partial x}\right) = \frac{1}{a^2}(a_y\,DW - a_x\,dZ)\exp[i(a_x\,x+a_y\,y)+\sigma t].
\end{align}
The foregoing equations relating $u$ and $v$ with the normal components
of the velocity and the vorticity are quite general: they are not restricted
to this particular problem.

\section{The principle of the exchange of stabilities}\label{sec:11}

We shall now show that, for the problem under discussion, the principle
of the exchange of stabilities is valid, i.e.\ $\sigma$ is real and the marginal states
are characterized by $\sigma = 0$.

Let
\begin{equation}
\label{eq:2-112}
G = (D^2 - a^2)W
\end{equation}
and
\begin{equation}
\label{eq:2-113}
F = (D^2-a^2)(D^2-a^2-\sigma)W = (D^2-a^2-\sigma)G.
\end{equation}
According to equation~(\ref{eq:2-99}), the condition $\Theta = 0$ on the bounding
surfaces is equivalent to
\begin{equation}
\label{eq:2-114}
F = 0 \qquad\text{for } z = 0 \text{ and } 1.
\end{equation}
In terms of $F$, the equation satisfied by $W$ is
\begin{equation}
\label{eq:2-115}
(D^2-a^2-\mathrm{p}\sigma)F = -Ra^2 W.
\end{equation}
We shall now prove that $\sigma$ is real for all positive $R$.

Multiply equation~(\ref{eq:2-115}) by $F^*$ (the complex conjugate of $F$) and
integrate over the range of $z$. We have
\begin{equation}
\label{eq:2-116}
\int_0^1 F^*(D^2-a^2-\mathrm{p}\sigma)F\,dz = -Ra^2\!\int_0^1 F^*W\,dz.
\end{equation}
By an integration by parts we obtain
\begin{equation}
\label{eq:2-117}
\int_0^1 F^*D^2F\,dz = -\int_0^1 |DF|^2\,dz,
\end{equation}
since the integrated part vanishes on account of the boundary condition~(\ref{eq:2-114}). Hence
\begin{equation}
\label{eq:2-118}
\int_0^1 \{|DF|^2 + (a^2+\mathrm{p}\sigma)|F|^2\}\,dz = Ra^2\!\int_0^1 WF^*\,dz.
\end{equation}
Consider the integral on the right-hand side. We have
\begin{equation}
\label{eq:2-119}
\int_0^1 WF^*\,dz = \int_0^1 W(D^2-a^2-\sigma^*)G^*\,dz.
\end{equation}
\begin{equation}
\label{eq:2-120}
= \int_0^1 WD^2G^*\,dz - (a^2+\sigma^*)\!\int_0^1 WG^*\,dz.
\end{equation}
By successive integrations by parts, we obtain
\begin{equation}
\label{eq:2-121}
\int_0^1 WD^2G^*\,dz = -\int_0^1 DWDG^*\,dz = \int_0^1 G^*D^2W\,dz,
\end{equation}
the integrated part vanishing each time on account of the boundary
conditions $W = 0$ and \emph{either} $DW = 0$ or $G^* = (D^2-a^2)W^* = 0$
(depending on whether the particular bounding surface is rigid or free).
Thus
\begin{equation}
\label{eq:2-122}
\int_0^1 WF^*\,dz = \int_0^1 G^*\{(D^2-a^2)W - \sigma^*W\}\,dz = \int_0^1 |G|^2\,dz - \sigma^*\!\int_0^1 WD^2W^*\,dz.
\end{equation}
Again, by an integration by parts, we have
\begin{equation}
\label{eq:2-123}
\int_0^1 WD^2W^*\,dz = -\int_0^1 |DW|^2\,dz.
\end{equation}
Combining equations~(\ref{eq:2-118}),~(\ref{eq:2-122}), and~(\ref{eq:2-123}), we obtain
\begin{equation}
\label{eq:2-124}
\int_0^1 \{|DF|^2+a^2|F|^2+\mathrm{p}\sigma|F|^2\}\,dz - Ra^2\!\int_0^1\{|G|^2+\sigma^*[|DW|^2+a^2|W|^2]\}\,dz = 0.
\end{equation}
The real and the imaginary parts of this equation must vanish separately;
and the vanishing of the imaginary part gives
\begin{equation}
\label{eq:2-125}
\operatorname{im}(\sigma)\!\left\{\mathrm{p}\!\int_0^1|F|^2\,dz + Ra^2\!\int_0^1[|DW|^2+a^2|W|^2]\,dz\right\} = 0.
\end{equation}
But the quantity inside the curly brackets is positive definite for $R > 0$.
Hence
\begin{equation}
\label{eq:2-126}
\operatorname{im}(\sigma) = 0.
\end{equation}
This establishes that $\sigma$ is real for $R > 0$ and that the principle of the
exchange of stabilities is valid for this problem.

\section{The equations governing the marginal state and the reduction to a characteristic value problem}\label{sec:12}

Since $\sigma$ is real for all positive Rayleigh numbers (i.e.\ for all adverse
temperature gradients), it follows that the transition from stability to
instability must occur via a stationary state. The equations governing
the marginal state are therefore to be obtained by setting $\sigma = 0$ in the
relevant equations. Thus, we have (cf.\ equations~(\ref{eq:2-99}) and~(\ref{eq:2-100}))
\begin{equation}
\label{eq:2-127}
(D^2-a^2)^2 W = +\left(\frac{g\alpha}{\kappa\nu}\,d^4\right)a^2\Theta
\end{equation}
and
\begin{equation}
\label{eq:2-128}
(D^2-a^2)\Theta = -\left(\frac{\beta}{\kappa}\,d^2\right)W.
\end{equation}
Eliminating $\Theta$ between these equations, we obtain
\begin{equation}
\label{eq:2-129}
(D^2-a^2)^3 W = -Ra^2 W.
\end{equation}
Solutions of this equation must be sought which satisfy the boundary
conditions
\begin{equation}
\label{eq:2-130}
W = 0, \quad (D^2-a^2)^2W = 0 \quad\text{for } z = 0 \text{ and } 1,
\end{equation}
and either $DW$ or $D^2W = 0$ depending on the nature of the
bounding surfaces at $z = 0$ and 1.

Alternatively, we can eliminate $W$ between the equations~(\ref{eq:2-127}) and~(\ref{eq:2-128}) and obtain
\begin{equation}
\label{eq:2-131}
(D^2-a^2)^3\Theta = -Ra^2\Theta,
\end{equation}
together with the boundary conditions
\begin{equation}
\label{eq:2-132}
\Theta = 0, \quad (D^2-a^2)\Theta = 0 \quad\text{for } z = 0 \text{ and } 1,
\end{equation}
and either $D(D^2-a^2)\Theta$ or $D\{D(D^2-a^2)\Theta\} = 0$ depending on the
nature of the bounding surfaces at $z = 0$ and~1.

In either of the two foregoing formulations (equations~(\ref{eq:2-129}) and~(\ref{eq:2-130}) or equations~(\ref{eq:2-131}) and~(\ref{eq:2-132})) we have a differential equation of order
six and we have to satisfy six boundary conditions, three at $z = 0$
and three at $z = 1$; this we cannot in general do. Only for particular
values of $R$ (for given $a^2$) will the problem allow a non-zero solution. We
have thus a \emph{characteristic value problem} for $R$.

The manner in which the solution of the stability problem is to be
carried out is clear. For a given $a^2$ we must determine the lowest characteristic value for $R$; the minimum with respect to $a^2$ of the characteristic
values so obtained is the critical Rayleigh number at which instability
will manifest itself.

\section{The variational principles}\label{sec:13}

We shall now show that the solution of the characteristic value
problems formulated in \S~12 can be expressed in terms of variational
principles. The first of these is due to Pellew and Southwell.

\subsection{The first variational principle}\label{sec:13a}

Consider first the characteristic value problem expressed in terms
of $W$. Letting
\begin{equation}
\label{eq:2-133}
F = (D^2-a^2)^2 W = (D^2-a^2)G,
\end{equation}
we have
\begin{equation}
\label{eq:2-134}
(D^2-a^2)F = -Ra^2 W;
\end{equation}
and the boundary conditions are:
\begin{equation}
\label{eq:2-135}
W = 0 \;\text{ and }\; F = 0 \;\text{ for } z = 0 \text{ and } 1; \;\text{ and \emph{either} } DW = 0 \;\text{ \emph{or} } D^2W = 0
\end{equation}
depending on the nature of the bounding surfaces at $z = 0$ and~1.

Let $R_1$ be a characteristic value; and let the solution belonging to $R_1$
be distinguished by a subscript $j$. Multiplying the equation satisfied by
$W_i$ by $F_j$ (belonging to a different characteristic value $R_j$) and integrating over the range of $z$, we have
\begin{equation}
\label{eq:2-136}
\int_0^1 F_i(D^2-a^2)F_j\,dz = -R_j\,a^2\!\int_0^1 W_i(D^2-a^2)G_j\,dz.
\end{equation}
By a sequence of integrations by parts similar to that followed in \S~11,
we obtain from equation~(\ref{eq:2-136}) the result (cf.\ equation~(\ref{eq:2-124}))
\begin{equation}
\label{eq:2-137}
\int_0^1(DF_i\,DF_j+a^2 F_i\,F_j)\,dz = R_j\,a^2\!\int_0^1 G_i\,G_j\,dz.
\end{equation}
Interchanging $i$ and $j$ in this last equation, we have
\begin{equation}
\label{eq:2-138}
\int_0^1(DF_j\,DF_i+a^2 F_j\,F_i)\,dz = R_i\,a^2\!\int_0^1 G_j\,G_i\,dz.
\end{equation}
From equations~(\ref{eq:2-137}) and~(\ref{eq:2-138}) it is apparent that
\begin{equation}
\label{eq:2-139}
\int_0^1 G_i\,G_j\,dz = 0 \qquad\text{if } i \neq j.
\end{equation}

\emph{The functions $G_j$ belonging to the different characteristic values are, therefore,
orthogonal.}

When $i = j$, equation~(\ref{eq:2-137}) gives
\begin{equation}
\label{eq:2-140}
R_j\,a^2\!\int_0^1 G_j^2\,dz = \int_0^1[(DF_j)^2+a^2 F_j^2]\,dz,
\end{equation}
which expresses $R_j$ as the ratio of two positive definite integrals. Indeed, we shall show that $R$ given by the formula
\begin{equation}
\label{eq:2-141}
R = \frac{\displaystyle\int_0^1[(DF)^2+a^2 F^2]\,dz}{a^2\!\displaystyle\int_0^1 G^2\,dz} = \frac{\mathscr{I}_1}{\mathscr{I}_2} \quad\text{(say)}
\end{equation}
has a stationary property when the quantities on the right-hand side
are evaluated in terms of the true characteristic functions.

To prove the stationary property, let $\mathscr{R}$ be evaluated in accordance
with equation~(\ref{eq:2-141}) in terms of a function $W$ which is arbitrary except
for the requirement (apart from boundedness and continuity) that it
satisfy the boundary conditions~(\ref{eq:2-135}). Let $\delta R$ be the change in $\mathscr{R}$ when
$W$ is subjected to a small variation $\delta W$ which is again compatible with the
boundary conditions on $W$; i.e.
\begin{equation}
\label{eq:2-142}
\delta W = 0 \text{ and } \delta F = 0 \text{ for } z = 0 \text{ and } 1;\text{ and either } D\,\delta W = 0 \text{ or } D^2\delta W = 0
\end{equation}
depending on the nature of the bounding surfaces at $z = 0$ and~1.

According to equation~(\ref{eq:2-141}),
\begin{equation}
\label{eq:2-143}
\delta R = \frac{1}{\mathscr{I}_2}(\delta\mathscr{I}_1 - \tfrac{\mathscr{I}_1}{\mathscr{I}_2}\delta\mathscr{I}_2) = \frac{1}{\mathscr{I}_2}(\delta\mathscr{I}_1 - Ra^2\delta\mathscr{I}_2),
\end{equation}
where
\begin{equation}
\label{eq:2-144}
\delta\mathscr{I}_1 = 2\!\int_0^1(DFD\delta F + a^2 F\,\delta F)\,dz
\end{equation}
and
\begin{equation}
\label{eq:2-145}
\delta\mathscr{I}_2 = 2\!\int_0^1[(D^2-a^2)W](D^2-a^2)\,\delta W\,dz
\end{equation}
are the variations in $\mathscr{I}_1$ and $\mathscr{I}_2$ consequent to the variation $\delta W$ in $W$. After
a sequence of integrations by parts, we readily find
\begin{equation}
\label{eq:2-146}
\delta\mathscr{I}_1 = -2\!\int_0^1 \delta F(D^2-a^2)F\,dz
\end{equation}
and
\begin{equation}
\label{eq:2-147}
\delta\mathscr{I}_2 = 2\!\int_0^1 W(D^2-a^2)^2\,\delta F\,dz = 2\!\int_0^1 W\,\delta F\,dz.
\end{equation}
Thus
\begin{equation}
\label{eq:2-148}
\delta R = -\frac{2}{\mathscr{I}_2}\!\int_0^1 \delta F\{(D^2-a^2)F + Ra^2 W\}\,dz.
\end{equation}

From~(\ref{eq:2-148}) it follows that $\delta R = 0$ \emph{if}
\begin{equation}
\label{eq:2-149}
(D^2-a^2)F = -Ra^2 W;
\end{equation}
\emph{and, conversely, if $\delta R = 0$ for any arbitrary $\delta F$ $[= (D^2-a^2)^2\delta W]$ compatible
with the boundary conditions of the problem, then equation}~(\ref{eq:2-149}) \emph{must hold
and the function $F$ in terms of which $R$ was initially calculated must have
been a solution of the characteristic value problem.}

The foregoing arguments establish the stationary property of the
characteristic values when they are regarded as given by equation~(\ref{eq:2-141}).
We shall now show, further, that the lowest characteristic value of $R$
is, indeed, a true minimum.

We have seen that the functions $G_j$ form an orthogonal set. We shall
suppose that they are also normalized so that
\begin{equation}
\label{eq:2-150}
\int_0^1 G_i\,G_j\,dz = \delta_{ij}.
\end{equation}
Let
\begin{equation}
\label{eq:2-151}
G = (D^2-a^2)W,
\end{equation}
where $W$ is an arbitrary, continuous, bounded function which satisfies
the boundary conditions of the problem. We shall suppose that $G$ can
be expanded in the basic set of functions $G_j$; thus,
\begin{equation}
\label{eq:2-152}
G = \sum_{j=1}^{\infty} A_j\,G_j,
\end{equation}
where
\begin{equation}
\label{eq:2-153}
A_j = \int_0^1 G\,G_j\,dz.
\end{equation}
Suppose $G$ is normalized: then
\begin{equation}
\label{eq:2-154}
1 = \int_0^1 G^2\,dz = \sum_j\sum_k A_j\,A_k\!\int_0^1 G_j\,G_k\,dz = \sum_j A_j^2.
\end{equation}

Associated with the expansion~(\ref{eq:2-152}) for $G_j$, we have
\begin{equation}
\label{eq:2-155}
W = \sum_{j=1}^{\infty} A_j\,W_j
\end{equation}
and
\begin{equation}
\label{eq:2-156}
F = \sum_{j=1}^{\infty} A_j(D^2-a^2)W_j = \sum_{j=1}^{\infty} A_j(D^2-a^2)G_j.
\end{equation}
From equation~(\ref{eq:2-156}) it follows that
\begin{equation}
\label{eq:2-157}
(D^2-a^2)F = \sum_{j=1}^{\infty} A_j(D^2-a^2)^2 W_j = -a^2\sum_{j=1}^{\infty} A_j\,R_j\!\int_0^1 W_j(D^2-a^2)W_k\,dz.
\end{equation}
Multiply this last equation by $F$ and integrate over the range of $z$. We
obtain
\begin{align}
\label{eq:2-158}
\int_0^1 F(D^2-a^2)F\,dz &= -a^2\sum_j\sum_k A_j\,R_j\!\left\{\sum_k A_k\!\int_0^1 W_j(D^2-a^2)W_k\,dz\right\} \notag\\
&= -a^2\sum_j\sum_k A_j\,A_k\!\int_0^1 (D^2-a^2)W_j(D^2-a^2)W_k\,dz \notag\\
&= -a^2\sum_j\sum_k A_j\,A_k\,R_j\!\int_0^1 G_j\,G_k\,dz \notag\\
&= -a^2\sum_{j=1}^{\infty} A_j^2\,R_j.
\end{align}
Thus
\begin{equation}
\label{eq:2-159}
\int_0^1[(DF)^2+a^2 F^2]\,dz = a^2\sum_{j=1}^{\infty} A_j^2\,R_j.
\end{equation}
or, making use of~(\ref{eq:2-154}), we can write
\begin{align}
\label{eq:2-160}
\int_0^1[(DF)^2+a^2 F^2]\,dz - a^2 R_1 &= a^2\!\left\{\sum_{j=1}^{\infty} A_j^2\,R_j - R_1\sum_{j=1}^{\infty} A_j^2\right\} \notag\\
&= a^2\sum_{j=2}^{\infty} A_j^2(R_j - R_1).
\end{align}
It is clear that the last summation in~(\ref{eq:2-160}) cannot be negative. Hence
\begin{equation}
\label{eq:2-161}
R_1\,a^2 \leqslant \int_0^1[(DF)^2+a^2 F^2]\,dz,
\end{equation}
where the equality sign can hold if and only if $A_j = 0$ for all $j$ greater
than 1. \emph{This proves that the quantity on the right-hand side of}~(\ref{eq:2-161}) \emph{attains
its true minimum when $F$ belongs to $R_1$.}

\subsection{The second variational principle}\label{sec:13b}

There is a second variational principle which one can obtain by
considering the characteristic value problem expressed in terms of $\Theta$.

Letting
\begin{equation}
\label{eq:2-162}
G' = (D^2-a^2)\Theta,
\end{equation}
we write the equation governing $\Theta$ in the form
\begin{equation}
\label{eq:2-163}
(D^2-a^2)^2 G' = -Ra^2\Theta;
\end{equation}
and the boundary conditions are
\begin{equation}
\label{eq:2-164}
\Theta = 0 \;\text{ and }\; G' = 0 \;\text{ for } z = 0 \text{ and } 1; \text{ and \emph{either} } DG' = 0 \;\text{ \emph{or} }\; D^2G' = 0
\end{equation}
depending on the nature of the bounding surfaces at $z = 0$ and~1.

Let the solution belonging to a particular characteristic value $R_j$ be
distinguished by a subscript $j$. Multiplying the equation satisfied by
$\Theta_j$ by $G'_i$ (belonging to $R_i$) and integrating over the range of $z$, we obtain
\begin{equation}
\label{eq:2-165}
\int_0^1 G'_i(D^2-a^2)^2 G'_j\,dz = -R_j\,a^2\!\int_0^1 \Theta_j(D^2-a^2)\Theta_i\,dz.
\end{equation}
By an integration by parts, the right-hand side becomes
\begin{equation}
\label{eq:2-166}
R_j\,a^2\!\int_0^1[(D\Theta_i)(D\Theta_j) + a^2\Theta_i\,\Theta_j]\,dz;
\end{equation}
the integrated parts vanish on account of the boundary condition $\Theta = 0$
for $z = 0$ and~1. Similarly, we obtain after successive integrations by
parts
\begin{align}
\label{eq:2-167}
\int_0^1 G'_i\,D^2(D^2-a^2)G'_j\,dz &= -\int_0^1 DG'_i\,D(D^2-a^2)G'_j\,dz \notag\\
&= +\int_0^1 D^2 G'_i\,(D^2-a^2)G'_j\,dz;
\end{align}
again the integrated parts vanish on account of the boundary conditions
on $G'$. Thus, we obtain
\begin{equation}
\label{eq:2-168}
R_j\,a^2\!\int_0^1[(D\Theta_i)(D\Theta_j)+a^2\Theta_i\,\Theta_j]\,dz = \int_0^1[(D^2-a^2)G'_i][(D^2-a^2)G'_j]\,dz.
\end{equation}
From~(\ref{eq:2-168}) it follows that
\begin{equation}
\label{eq:2-169}
\int_0^1(D^2-a^2)G'_i\,[(D^2-a^2)G'_j]\,dz = 0 \qquad\text{if } i \neq j;
\end{equation}
and that (when $i = j$)
\begin{equation}
\label{eq:2-170}
R_j = \frac{\displaystyle\int_0^1[(D^2-a^2)G'_j]^2\,dz}{a^2\!\displaystyle\int_0^1[(D\Theta_j)^2+a^2\Theta_j^2]\,dz}.
\end{equation}

Again it can be shown that $R$ given by the foregoing formula has a
stationary property and that the true minimum of the quantity on the
right-hand side is the lowest characteristic value of $R$.

\section{The thermodynamic significance of the variational principle}\label{sec:14}

We have seen that the Rayleigh number, at which disturbances of an
assigned wave number become unstable, is the minimum value which a
certain ratio of two positive definite integrals can attain. We shall now
show that the quantity which is thus minimized has a simple physical
meaning.

First we observe that when marginal conditions prevail, the $z$-component of the vorticity vanishes. This follows from equation~(\ref{eq:2-94}); for,
when $\sigma = 0$,
\begin{equation}
\label{eq:2-171}
(D^2-a^2)Z = 0;
\end{equation}
and this equation does not allow a non-zero solution which satisfies the
boundary conditions (cf.\ equation~(\ref{eq:2-96})) on $Z$. [The vanishing of $Z$ under
stationary conditions follows even more generally from equation~(\ref{eq:2-73});
for, when $\partial\zeta/\partial t = 0$, $\nabla^2\zeta = 0$ and as is well known, Laplace's equation
does not allow non-zero solutions which (or, the normal derivatives of
which) vanish on a closed boundary.] Hence, in this case, the solutions
for the horizontal components of the velocity given by equations~(\ref{eq:2-110})
and~(\ref{eq:2-111}) become
\begin{equation}
\label{eq:2-172}
u = \frac{1}{a^2}\frac{\partial^2 w}{\partial x\partial z} \quad\text{and}\quad v = \frac{1}{a^2}\frac{\partial^2 w}{\partial y\partial z}.
\end{equation}

To be specific, let
\begin{equation}
\label{eq:2-173}
w = W(z)\cos a_x\,x\cos a_y\,y,
\end{equation}
where
\begin{equation}
\label{eq:2-174}
a^2 = a_x^2+a_y^2.
\end{equation}
Then
\begin{equation}
\label{eq:2-175}
u = -\frac{DW}{a^2}\,a_x\sin a_x\,x\cos a_y\,y; \qquad v = -\frac{DW}{a^2}\,a_y\cos a_x\,x\sin a_y\,y.
\end{equation}

Now consider the average rate of viscous dissipation of energy by a
unit (vertical) column of the fluid. This is given by\footnote{Note that in this equation (as in all equations after~(\ref{eq:2-98})) length is measured in the unit $d$; this is the origin of the factor $\frac{1}{d^2}$ in this equation.}
\begin{align}
\label{eq:2-176}
\epsilon_v &= -\frac{\rho\nu}{d^2}\!\int_0^1\{(w\nabla^2 w)+(u\nabla^2 u)+(v\nabla^2 v)\}\,dz \notag\\
&= -\frac{\rho\nu}{d^2}\!\int_0^1\{(w(D^2-a^2)w)+(u(D^2-a^2)u)+(v(D^2-a^2)v)\}\,dz,
\end{align}
where angular brackets signify that the quantity enclosed is averaged
over the horizontal plane. For $w$, $u$, and $v$ given by equations~(\ref{eq:2-173}) and~(\ref{eq:2-175}), we have
\begin{align}
\label{eq:2-177}
\epsilon_v &= -\frac{\rho\nu}{4d^2}\!\int_0^1\!\left\{W(D^2-a^2)W + \frac{a_x^2}{a^2}\overline{DW}(D^2-a^2)DW + \right. \notag\\
&\qquad\qquad\qquad\left. +\frac{a_y^2}{a^2}DW(D^2-a^2)DW\right\}dz \notag\\
&= -\frac{\rho\nu}{4d^2}\!\int_0^1\{a^2 W(D^2-a^2)W + DW(D^2-a^2)DW\}\,dz.
\end{align}
Integrating by parts the second term on the right-hand side, we obtain
(cf.\ equation~(\ref{eq:2-133}))
\begin{equation}
\label{eq:2-178}
\epsilon_v = \frac{\rho\nu}{4d^2}\!\int_0^1[(D^2-a^2)W]^2\,dz = \frac{\rho\nu}{4d^2}\!\int_0^1 G^2\,dz.
\end{equation}

Consider next the average rate at which energy is released in a unit
column of the fluid by the buoyancy force $g\,\delta\rho$ $(= -g\alpha\rho\theta)$ acting on the
fluid. This is given by
\begin{equation}
\label{eq:2-179}
\epsilon_g = \rho g\alpha\!\int_0^1(\theta w)\,dz.
\end{equation}
Making use of equation~(\ref{eq:2-76}) (and remembering that in this equation
length is \emph{not} measured in the unit $d$) we can rewrite~(\ref{eq:2-179}) in the form
\begin{equation}
\label{eq:2-180}
\epsilon_g = -\frac{\rho g\alpha\kappa}{4\beta d^2}\!\int_0^1\Theta(D^2-a^2)\Theta\,dz.
\end{equation}
After an integration by parts, we have
\begin{equation}
\label{eq:2-181}
\epsilon_g = \frac{\rho g\alpha\kappa}{4\beta d^2}\!\int_0^1[(D\Theta)^2+a^2\Theta^2]\,dz.
\end{equation}
On the other hand, according to equations~(\ref{eq:2-127}) and~(\ref{eq:2-133})
\begin{equation}
\label{eq:2-182}
\Theta = \frac{\nu}{g\alpha d^2}(D^2-a^2)^2 W = \frac{\nu}{g\alpha d^2}F.
\end{equation}
In terms of $F$, the expression for $\epsilon_g$ becomes
\begin{equation}
\label{eq:2-183}
\epsilon_g = \frac{\rho\nu^2\kappa}{4g\alpha\beta d^6}\!\int_0^1[(DF)^2+a^2 F^2]\,dz.
\end{equation}

In a steady state, the kinetic energy dissipated by viscosity must be
balanced by the internal energy released by the buoyancy force (see
Appendix~I). We must therefore require that
\begin{equation}
\label{eq:2-184}
\epsilon_v = \epsilon_g.
\end{equation}
Equations~(\ref{eq:2-178}) and~(\ref{eq:2-183}) now give
\begin{equation}
\label{eq:2-185}
\frac{\rho\nu}{4d^2}\!\int_0^1 G^2\,dz = \frac{\rho\nu^2\kappa}{4g\alpha\beta d^6}\!\int_0^1[(DF)^2+a^2 F^2]\,dz;
\end{equation}
or, alternatively,
\begin{equation}
\label{eq:2-186}
R = \frac{g\alpha\beta}{\kappa\nu}\,d^4 = \frac{\displaystyle\int_0^1[(DF)^2+a^2 F^2]\,dz}{a^2\!\displaystyle\int_0^1 G^2\,dz}.
\end{equation}
This is the same as equation~(\ref{eq:2-141}) which led to the first of the two
variational principles considered in \S~13. We may therefore state the
physical content of this variational principle as follows.

\emph{Instability occurs at the minimum temperature gradient at which a
balance can be steadily maintained between the kinetic energy dissipated
by viscosity and the internal energy released by the buoyancy force.}

\section{Exact solutions of the characteristic value problem}\label{sec:15}

We now return to the characteristic value problem formulated in
\S~12. We shall obtain the solution for the following three cases.

\begin{enumerate}[label=(\alph*)]
\item Both the bounding surfaces are free, in which case
\begin{equation}
\label{eq:2-187}
W = (D^2-a^2)^2 W = 0 \quad\text{and}\quad D^2 W = 0 \;\text{ for } z = 0 \text{ and } 1.
\end{equation}
\item Both the bounding surfaces are rigid, in which case
\begin{equation}
\label{eq:2-188}
W = (D^2-a^2)^2 W = 0 \quad\text{and}\quad DW = 0 \;\text{ for } z = 0 \text{ and } 1.
\end{equation}
\item One of the bounding surfaces (say, $z = 0$) is rigid and the other
bounding surface ($z = 1$) is free, in which case
\begin{equation}
\label{eq:2-189}
W = (D^2-a^2)^2 W = 0 \;\text{ for } z = 0 \text{ and } 1; \text{ and } DW = 0 \;\text{ for } z = 0 \;\text{ and}\;
D^2 W = 0 \;\text{ for } z = 1.
\end{equation}
\end{enumerate}

From the point of view of realizability in the laboratory for precise
quantitative experiments, case~(b) is, of course, of the greatest interest.
Case~(c) is of interest when the top surface is kept free for visual observations (as in B\'enard's original experiments). Case~(a) (which was first
considered by Rayleigh) can be realized, if at all, only under very
artificial conditions (e.g.\ by having the liquid layer floating on top of a
somewhat heavier liquid); nevertheless, the case is of theoretical interest
because it allows an explicit solution and clearly this has some
advantages.

\subsection{The solution for two free boundaries}\label{sec:15a}

In this case the boundary conditions~(\ref{eq:2-187}) require
\begin{equation}
\label{eq:2-190}
W = D^2W = D^4W = 0 \qquad\text{for } z = 0 \text{ and } 1;
\end{equation}
from the equation satisfied by $W$ (namely~(\ref{eq:2-129})) it now follows that
$D^6W = 0$ for $z = 0$ and~1. From equation~(\ref{eq:2-129}) differentiated twice
with respect to $z$, we next conclude that $D^8W = 0$ is also zero for $z = 0$
and~1. By further differentiations of equation~(\ref{eq:2-129}), we can conclude,
successively, that all the even derivatives of $W$ vanish on the boundaries.
Thus
\begin{equation}
\label{eq:2-191}
D^{2m}W = 0 \qquad\text{for } z = 0 \text{ and } 1 \quad\text{and}\quad m = 1, 2, \ldots\,.
\end{equation}
From this it follows that the required solutions must be
\begin{equation}
\label{eq:2-192}
W = A\sin n\pi z \qquad(n = 1, 2, \ldots),
\end{equation}
where $A$ is a constant and $n$ is an integer. Substitution of this solution
in equation~(\ref{eq:2-129}) leads to the \emph{characteristic equation}
\begin{equation}
\label{eq:2-193}
R = (n^2\pi^2+a^2)^3/a^2.
\end{equation}
For a given $a^2$, the lowest value of $R$ occurs when $n = 1$; then
\begin{equation}
\label{eq:2-194}
R = \frac{(\pi^2+a^2)^3}{a^2}.
\end{equation}
The meaning of this last relation is this: for all Rayleigh numbers less
than that given by~(\ref{eq:2-194}), disturbances with the wave number $a$ will be
stable; these disturbances will become marginally stable when the
Rayleigh number equals the value given by~(\ref{eq:2-194}); and when the Rayleigh
number exceeds the value given by~(\ref{eq:2-194}), the same disturbances will
be unstable. The critical Rayleigh number for the onset of instability
is, therefore, determined by the condition
\begin{equation}
\label{eq:2-195}
\frac{\partial R}{\partial a^2} = 3\frac{(\pi^2+a^2)^2}{a^2}-\frac{(\pi^2+a^2)^3}{a^4} = 0,
\end{equation}
or, \qquad $3a^2 = \pi^2+a^2$ \quad or \quad $a^2 = \pi^2/2$.

The corresponding value of $R$ is
\begin{equation}
\label{eq:2-196}
R_c = \frac{(\frac{3}{2}\pi^2)^3}{\frac{1}{2}\pi^2} = \frac{27}{4}\pi^4 = 657\!\cdot\!5;
\end{equation}
and the disturbances which will be manifest at marginal stability will
be characterized by the wavelength
\begin{equation}
\label{eq:2-197}
\lambda = \frac{2\pi}{k} = \frac{2\pi}{a}\,d = 2\sqrt{2}\,d.
\end{equation}

\subsection{The solution for two rigid boundaries}\label{sec:15b}

In view of the symmetry of the two bounding planes with respect to
the two bounding planes, it will be convenient to translate the origin of $z$ to be
midway between the two planes. Then, the fluid will be confined between
$z = \pm\frac{1}{2}$; and we have to seek solutions of the equation
\begin{equation}
\label{eq:2-198}
(D^2-a^2)^3 W = -Ra^2 W
\end{equation}
which satisfy the boundary conditions
\begin{equation}
\label{eq:2-199}
W = DW = (D^2-a^2)^2 W = 0 \qquad\text{for } z = \pm\tfrac{1}{2}.
\end{equation}

We may first observe that it follows from the evenness of the operator
$(D^2-a^2)^3$, and the identity of the boundary conditions which have to be
satisfied at $z = \pm\frac{1}{2}$, that the proper solutions of equation~(\ref{eq:2-198}) fall
into two non-combining groups of even and odd solutions. From general
considerations, it is clear that the lowest `state' will be even with no nodes
while the first `excited state' will be odd with a node at $z = 0$.

It is evident that the general solution of equation~(\ref{eq:2-198}) can be expressed as a superposition of solutions of the form
\begin{equation}
\label{eq:2-200}
W = e^{qz},
\end{equation}
where $q^2$ is a root of the equation
\begin{equation}
\label{eq:2-201}
(q^2-a^2)^3 = -Ra^2.
\end{equation}
Letting
\begin{equation}
\label{eq:2-202}
Ra^2 = \tau^3 a^6,
\end{equation}
we find that the roots of equation~(\ref{eq:2-201}) are given by
\begin{equation}
\label{eq:2-203}
q^2 = -a^2(\tau-1) \quad\text{and}\quad q^2 = a^2[1+\tfrac{1}{2}\tau(1\pm i\sqrt{3})];
\end{equation}
or, taking the square roots, we have the six roots\footnote{An asterisk to a quantity denotes its complex conjugate.}
\begin{equation}
\label{eq:2-204}
\pm iq_0, \quad \pm q, \quad\text{and}\quad \pm q^*,
\end{equation}
where
\begin{equation}
\label{eq:2-205}
q_0 = a(\tau - 1)^{1/2},
\end{equation}
\begin{align}
\label{eq:2-206}
\operatorname{re}(q) &= q_1 = a\{\tfrac{1}{2}[\sqrt{1+\tau+\tau^2}+\tfrac{1}{2}(1+\tfrac{1}{2}\tau)]\}^{1/2}, \notag\\
\operatorname{im}(q) &= q_2 = a\{\tfrac{1}{2}[\sqrt{1+\tau+\tau^2}-\tfrac{1}{2}(1+\tfrac{1}{2}\tau)]\}^{1/2}.
\end{align}
From~(\ref{eq:2-203}) we readily obtain the following relations which we shall
find useful:
\begin{equation}
\label{eq:2-207}
(q_0^2+a^2)^2 = a^4\tau^2
\end{equation}
and
\begin{equation}
\label{eq:2-208}
(q^2-a^2)^2 = \tfrac{1}{4}a^4\tau^2(-1\pm i\sqrt{3}).
\end{equation}

\paragraph{(i) The even solutions}

Considering first the even solutions, we can clearly write
\begin{equation}
\label{eq:2-209}
W = A_0\cos q_0\,z + A\cosh qz + A^*\cosh q^*z,
\end{equation}
where $A_0$ and $A$ are constants, the latter being complex. We must now
impose the boundary conditions~(\ref{eq:2-199}) on this solution.

From~(\ref{eq:2-209}) we obtain
\begin{equation}
\label{eq:2-210}
DW = -A_0\,q_0\sin q_0\,z + Aq\sinh qz + A^*q^*\sinh q^*z,
\end{equation}
and
\begin{equation}
\label{eq:2-211}
(D^2-a^2)^2W = A_0(q_0^2+a^2)^2\cos q_0\,z + A(q^2-a^2)^2\cosh qz + A^*(q^{*2}-a^2)^2\cosh q^*z.
\end{equation}
Using the relations~(\ref{eq:2-207}), we can rewrite~(\ref{eq:2-211}) in the form
\begin{equation}
\label{eq:2-212}
(D^2-a^2)^2W = \tfrac{1}{4}a^4\tau^2[2A_0\cos q_0\,z + A(i\sqrt{3}-1)\cosh qz - (i\sqrt{3}+1)A^*\cosh q^*z].
\end{equation}
The boundary conditions~(\ref{eq:2-199}), therefore, require
\begin{equation}
\label{eq:2-213}
\begin{vmatrix}
\cos\frac{1}{2}q_0 & \cosh\frac{1}{2}q & \cosh\frac{1}{2}q^* \\
-q_0\sin\frac{1}{2}q_0 & q\sinh\frac{1}{2}q & q^*\sinh\frac{1}{2}q^* \\
\cos\frac{1}{2}q_0 & \frac{1}{2}(i\sqrt{3}-1)\cosh\frac{1}{2}q & -\frac{1}{2}(i\sqrt{3}+1)\cosh\frac{1}{2}q^*
\end{vmatrix} = 0.
\end{equation}
For a non-vanishing solution, the determinant of the matrix in~(\ref{eq:2-213}) must vanish. This condition is
\begin{equation}
\label{eq:2-214}
\begin{vmatrix}
1 & 1 & 1 \\
-q_0\tan\frac{1}{2}q_0 & q\tanh\frac{1}{2}q & q^*\tanh\frac{1}{2}q^* \\
1 & \frac{1}{2}(i\sqrt{3}-1) & -\frac{1}{2}(i\sqrt{3}+1)
\end{vmatrix} = 0.
\end{equation}
Subtracting the first row from the third row and dividing the result by
$-\sqrt{3}/2$, we get
\begin{equation}
\label{eq:2-215}
\begin{vmatrix}
1 & 1 & 1 \\
-q_0\tan\frac{1}{2}q_0 & q\tanh\frac{1}{2}q & q^*\tanh\frac{1}{2}q^* \\
0 & \sqrt{3}-i & \sqrt{3}+i
\end{vmatrix} = 0.
\end{equation}
Expanding this last determinant, we have
\begin{equation}
\label{eq:2-216}
\operatorname{im}\{(\sqrt{3}+i)(q_1+iq_2)\tanh\frac{1}{2}q\} + q_0\tan\frac{1}{2}q_0 = 0.
\end{equation}

\begin{table}[htbp]
\centering
\caption{The exact characteristic values for the first even and odd modes of instability}
\label{tab:2-I}
\medskip
\begin{tabular}{@{}c cc cc@{}}
\toprule
& \multicolumn{2}{c}{Even solution} & \multicolumn{2}{c}{Odd solution} \\
\cmidrule(lr){2-3} \cmidrule(lr){4-5}
$a$ & $R$ & $a(\tau-1)^{\frac{1}{2}}$ & $R$ & $a(\tau-1)^{\frac{1}{2}}$ \\
\midrule
$0^+$    & $\infty$  & $\pi$       & $\infty$  & $7{\cdot}3313$   \\
$1{\cdot}0$ & $5854{\cdot}8$ & $4{\cdot}143514$ & $15312{\cdot}6$ & $7{\cdot}333940$  \\
$1{\cdot}5$ & $2737{\cdot}1$ & $4{\cdot}071304$ & $4700{\cdot}5$ & $7{\cdot}059910$  \\
$2{\cdot}0$ & $2171{\cdot}28$ & $3{\cdot}949002$ & $4544{\cdot}6$ & $7{\cdot}052095$  \\
$3{\cdot}117$ & $1797{\cdot}762$ & $3{\cdot}923639$ & $2498{\cdot}2$ & $7{\cdot}3569$ \\
$4{\cdot}0$ & $1859{\cdot}10$ & $3{\cdot}985334$ & $1068{\cdot}6$ & $7{\cdot}143707$ \\
$4{\cdot}0$ & $2439{\cdot}23$ & $3{\cdot}786334$ & $1773{\cdot}1$ & $7{\cdot}158372$  \\
$5{\cdot}365$ & $3417{\cdot}98$ & $3{\cdot}706519$ & $1761{\cdot}39$ & $7{\cdot}143877$  \\
$6{\cdot}0$ & $4018{\cdot}54$ & $3{\cdot}653524$ & $1953{\cdot}8$ & $7{\cdot}044460$  \\
$8{\cdot}0$ & $7084{\cdot}51$ & $3{\cdot}581053$ & $2240{\cdot}5$ & $6{\cdot}989981$ \\
 & & & & $6{\cdot}833135$ \\
\bottomrule
\end{tabular}
\end{table}

We can write this equation, alternatively, as
\begin{equation}
\label{eq:2-217}
-q_0\tan\frac{1}{2}q_0 = \operatorname{im}\!\left\{(\sqrt{3}+i)(q_1+iq_2)\frac{\sinh q_1+i\sin q_2}{\cosh q_1+\cos q_2}\right\}\!.
\end{equation}
Simplifying the right-hand side of equation~(\ref{eq:2-217}), we have
\begin{equation}
\label{eq:2-218}
-q_0\tan\frac{1}{2}q_0 = \frac{(q_1+q_2\sqrt{3})\sinh q_1+(q_1\sqrt{3}-q_2)\sin q_2}{\cosh q_1+\cos q_2}.
\end{equation}

Equation~(\ref{eq:2-218}) is a transcendental equation relating $a$ and
$\tau$ $(= \sqrt[3]{R/a^2})$, and it must be solved numerically by trial and error.
The method is to determine $\tau$ for a given $a$; the corresponding characteristic
value of $R$ then follows from equation~(\ref{eq:2-202}). In this way the solution
of this problem was first accomplished by Pellew and Southwell, though
the exact solution had been obtained earlier by a somewhat different
(but equivalent) method by Low. However, the definitive treatment of
the problem is due to Reid and Harris; their results are given in Table~I
and illustrated in Fig.~2. It will be observed that the function attains
its minimum at
\begin{equation}
\label{eq:2-219}
a = 3{\cdot}117, \quad\text{where } R = 1707{\cdot}762.
\end{equation}
The corresponding solutions for $W$ and $F$ $(= (a^2R)^{-1}F$ in~(\ref{eq:2-209}) set
equal to~1):
\begin{align}
\label{eq:2-220}
W &= \cos q_0\,z - 0{\cdot}06151664\cosh q_1\,z\cos q_2\,z + 0{\cdot}10388700\sinh q_1\,z\sin q_2\,z, \notag\\
(a^2R)^{-1}F &= \cos q_0\,z + 0{\cdot}12072710\cosh q_1\,z\cos q_2\,z + {} \notag\\
&\qquad\qquad\qquad{} + 0{\cdot}001331473\sinh q_1\,z\sin q_2\,z.
\end{align}

\figurePlaceholder{2}{The Rayleigh numbers at which instability sets in for disturbances of different wave numbers $a$ for the first even (curve labelled 1) and the first odd (curve labelled 2) mode.}

where
\begin{equation}
\label{eq:2-221}
q_0 = 3{\cdot}973839; \quad q_1 = 5{\cdot}195214; \quad q_2 = 2{\cdot}126096.
\end{equation}
The solutions $W$ and $F$ are given in Table~II and illustrated in Fig.~3.

\paragraph{(ii) The odd solutions}

Considering next the odd solutions, we now have (instead of~(\ref{eq:2-209}))
\begin{equation}
\label{eq:2-222}
W = A_0\sin q_0\,z + A\sinh qz + A^*\sinh q^*z.
\end{equation}
The characteristic determinant which follows from this solution is
(cf.\ equation~(\ref{eq:2-215}))
\begin{equation}
\label{eq:2-223}
\begin{vmatrix}
1 & 1 & 1 \\
q_0\cot\frac{1}{2}q_0 & q\coth\frac{1}{2}q & q^*\coth\frac{1}{2}q^* \\
0 & \sqrt{3}-i & \sqrt{3}+i
\end{vmatrix} = 0,
\end{equation}
and this leads to the equation
\begin{equation}
\label{eq:2-224}
q_0\cot\frac{1}{2}q_0 = \operatorname{im}\!\left\{(\sqrt{3}+i)(q_1+iq_2)\frac{\sinh q_1-i\sin q_2}{\cosh q_1-\cos q_2}\right\}\!.
\end{equation}
On simplifying this last equation, we obtain
\begin{equation}
\label{eq:2-225}
q_0\cot\frac{1}{2}q_0 = \frac{(q_1+q_2\sqrt{3})\sinh q_1-(q_1\sqrt{3}-q_2)\sin q_2}{\cosh q_1-\cos q_2}.
\end{equation}

The results of Reid and Harris derived on the basis of the solution~(\ref{eq:2-225}) are included in Tables~I and~II; and the solutions $W$ and $F$ for
the lowest odd mode are (see also Fig.~3):
\begin{align}
\label{eq:2-226}
W &= \sin q_0\,z - 0{\cdot}01707389\sinh q_1\,z\cos q_2\,z + 0{\cdot}00345645\cosh q_1\,z\sin q_2\,z, \notag\\
(a^2R)^{-1}F &= \sin q_0\,z + 0{\cdot}01153032\sinh q_1\,z\cos q_2\,z + {} \notag\\
&\qquad\qquad\qquad{} + 0{\cdot}01305820\cosh q_1\,z\sin q_2\,z,
\end{align}
where
\begin{equation}
\label{eq:2-227}
q_0 = 7{\cdot}137877; \quad q_1 = 9{\cdot}110819; \quad q_2 = 3{\cdot}789330.
\end{equation}
The first excited mode for the value of $a$ $(= 3{\cdot}117)$ at which the lowest
even mode gives its minimum has $R = 17\,610{\cdot}39$.

\begin{table}[htbp]
\centering
\caption{The solutions for $W$ and $(a^2 R)^{-1}F$ for the first even and odd modes of instability}
\label{tab:2-II}
\medskip
\begin{tabular}{@{}c cccc@{}}
\toprule
$z$ & $W_{\text{even}}$ & $(a^2R)^{-1}F_{\text{even}}$ & $W_{\text{odd}}$ & $(a^2R)^{-1}F_{\text{odd}}$ \\
\midrule
$0{\cdot}00$ & $+0{\cdot}938434$ & $+1{\cdot}120727$ & $0{\cdot}000000$ & $0{\cdot}000000$\\
$0{\cdot}05$ & $+0{\cdot}937350$ & $+1{\cdot}100728$ & $+0{\cdot}069030$ & $+0{\cdot}072860$\\
$0{\cdot}10$ & $+0{\cdot}931100$ & $+1{\cdot}181891$ & $+0{\cdot}130497$ & $+0{\cdot}145830$\\
$0{\cdot}15$ & $+0{\cdot}915004$ & $+1{\cdot}148709$ & $+0{\cdot}203157$ & $+0{\cdot}207189$\\
$0{\cdot}20$ & $+0{\cdot}890138$ & $+1{\cdot}110353$ & $+0{\cdot}274005$ & $+0{\cdot}288004$\\
$0{\cdot}25$ & $+0{\cdot}792804$ & $+1{\cdot}014961$ & $+0{\cdot}348796$ & $+0{\cdot}374315$\\
$0{\cdot}30$ & $+0{\cdot}709730$ & $+0{\cdot}917372$ & $+0{\cdot}406450$ & $+0{\cdot}453906$\\
$0{\cdot}35$ & $+0{\cdot}490256$ & $+0{\cdot}707835$ & $+0{\cdot}348446$ & $+0{\cdot}455965$\\
$0{\cdot}40$ & $+0{\cdot}118970$ & $+0{\cdot}409438$ & $+0{\cdot}188600$ & $+0{\cdot}343300$\\
$0{\cdot}45$ & $-0{\cdot}443530$ & $+0{\cdot}009139$ & $-0{\cdot}028960$ & $+0{\cdot}080456$\\
$0{\cdot}50$ & $0{\cdot}000000$ & $0{\cdot}000000$ & $-0{\cdot}300000$ & $-0{\cdot}300000$\\
\bottomrule
\end{tabular}
\end{table}

\figurePlaceholder{3}{The proper solutions for $W$ (curve 1) and $(a^2 R)^{-1}F$ (curve 2) for the state of marginal stability for the case when both bounding surfaces are rigid.}

